\textbf{Duración:} 3 horas (en aula)\\
\textbf{Enfoque:} Del brief a la arquitectura pedagógica de una obra técnica y un outline alineado a resultados de aprendizaje.

\subsubsection*{Objetivos de la sesión}
Al finalizar, la persona participante:
\begin{enumerate}
\item Traduce un brief en arquitectura de obra (capítulos, RA, artefactos).
\item Produce un outline de 6--10 nodos con learning outcomes en nodos clave.
\item Bosqueja un mapa preliminar de notación/glosario para consistencia STEM.
\item Justifica el alcance (qué queda fuera) para un microartefacto viable.
\end{enumerate}

\subsubsection*{Agenda (resumen)}
\begin{tabularx}{\textwidth}{@{}lclX@{}}
\toprule
Bloque & Tiempo & Contenido \\
\midrule
Recap & 10 min & Brief/orquestación $\rightarrow$ estructura \\
Arquitectura STEM & 30 min & Capítulos, RA, notación, ejercicios \\
Demo outline & 20 min & Brief $\rightarrow$ outline con restricciones \\
Taller & 70 min & Arquitectura + outline + mapa de símbolos \\
Pares & 25 min & Coherencia pedagógica \\
Cierre & 15 min & TI avance microartefacto \\
\bottomrule
\end{tabularx}

\subsubsection*{Producto esperado}
Arquitectura de obra + outline documentado con RA (E3).

\subsubsection*{Trabajo independiente relacionado}
Sesión 4: secciones, ejercicios y consistencia. TI: avanzar outline hacia primera sección (bloque C).
